\setcounter{chapter}{2}
\chapter{bootstrap}
\setcounter{section}{2}
\section{bootstrapによる推測}
\paragraph{表3.3}
最右列の理論値を求める．

まず記号の整理を行う．
求めたい指標は中央値：
\[ g(F) = \theta ; \quad g(F_n) = \tilde\theta . \]
bootstrapで求めるのは平均値：
\[ \hat\theta ; \quad \hat{\theta}^\ast . \]
バイアスは
\[ \text{Bias} = E(\hat\theta) - \theta , \quad \text{Bias}^\ast = E^\ast(\hat{\theta}^\ast) - \tilde\theta . \]

\(\hat\theta = \hat\theta(x_1 , \ldots , x_n)\)は平均値であり，今回の場合は
\[ \hat\theta = \frac{1}{n} \sum_{i=1}^{n} x_i = 218.9 . \]
bootstrapの対応する統計量は
\[ \hat{\theta}^\ast = \frac{1}{n} \sum_{i=1}^{n} x_i{}^\ast . \]
\(\tilde\theta = g(F_n)\)は（\(x_1, \ldots, x_n\)から得られる）経験的分布\(F_n\)の中央値であり，
\[ \tilde\theta = 129 . \]
さらに
\[
  E^\ast(\hat{\theta}^\ast)
  = \frac{1}{n} \sum_{i=1}^{n} E^\ast(x_i{}^\ast)
  = \frac{1}{n} \sum_{i=1}^{n} \frac{x_1 + \cdots + x_n}{n}
  = \frac{1}{n} \sum_{i=1}^{n} x_i
  = 218.9 .
\]
なので
\[ \text{Bias}^\ast = E^\ast(\hat{\theta}^\ast) - \tilde\theta = 218.9 - 129 = 89.9 . \]
したがって，\(\theta\)の推定値\(\hat\theta_c\)は
\[ \hat\theta_c = \hat\theta - \text{Bias}^\ast = 218.9 - 89.9 = 129 . \]

最後に\(\hat{\theta}^\ast\)の分散を求める．
\begin{align*}
  \Var^\ast(\hat{\theta}^\ast) &= E^\ast((\hat{\theta}^\ast)^2) - (E^\ast(\hat{\theta}^\ast))^2 \\
  &= E^\ast\left( \frac{1}{n^2} \sum_{ij} x_i{}^\ast x_j{}^\ast \right) - \left( \frac{1}{n} \sum_i x_i \right)^2 \\
  &= \frac{1}{n} \sum_i E^\ast((x_i{}^\ast)^2)
  + \frac{1}{n^2} \sum_{i\neq j} E^\ast(x_i{}^\ast x_j{}^\ast)
  - \frac{1}{n^2} \left( \sum_i x_i \right)^2 \\
  %
  &= \frac{1}{n} \sum_i \frac{x_1{}^2 + \cdots + x_n{}^2}{n}
  + \frac{1}{n^2} \sum_{j\neq j} \frac{(x_1 + \cdots + x_n)^2}{n^2}
  - \frac{1}{n^2} \left( \sum_i x_i \right)^2 \\
  %
  &= \frac{1}{n^2} \sum_i x_i{}^2 - \frac{1}{n^3} \left( \sum_i x_i \right)^2 \\
  &= 5330 .
\end{align*}

\section*{練習問題}
\addcontentsline{toc}{section}{練習問題}
\paragraph{問題3.3}
\(F(x)\)からの\(n\)個のサンプルを\(X_1, \ldots, X_n\)，及びこれらの順序統計量を\(Y_1, \ldots, Y_n\)とする．
\(Y_k\)の確率密度函数は
\[ g_k(y_k) = \frac{n!}{(k-1)!(n-k)!} [F(y_k)]^{k-1} [1 - F(y_k)]^{n-k} f(y_k) . \]
十分大きな\(n\)に対して\(n = kp\)としたときの\(Y_{np}\)の分散を求めたい．

\(\delta = F(y_k) - p\)で新しい確率変数を定義する．\(\delta\)の確率密度函数は
\begin{align*}
  h(\delta) &= g(y_k) \frac{dy_k}{d\delta} \\
  &= \frac{n!}{(k-1)!(n-k)!} [F(y_k)]^{k-1} [1 - F(y_k)]^{n-k} \\
  &= \frac{n!}{(np-1)!(n-np)!} [\delta + p]^{np-1} [1 - \delta - p]^{n-np} \\
  &\approx \frac{n!}{(np)!(n-np)!} [\delta + p]^{np} [1 - \delta - p]^{n(1-p)} \\
\end{align*}
対数をとれば
\begin{align*}
  \log h(\delta) &= \log \frac{n!}{(np)!(n-np)!} \left[ np \log(p + \delta) + n(1-p) \log(1-p-\delta) \right] \\
  &\approx \log \frac{n!}{(np)!(n-np)!} \left[ np \left( \log p + \frac{\delta}{p} - \frac{\delta^2}{2p^2} \right) + n(1-p) \left( \log(1-p) + \frac{-\delta}{1-p} - \frac{\delta^2}{2(1-p)^2} \right) \right] \\
  &= \log \frac{n!}{(np)!(n-np)!} \left[ np \log p + n(1-p) \log (1-p) - \frac{n}{p(1-p)} \frac{\delta^2}{2} \right] .
\end{align*}
Stirlingの公式\(n! \approx \sqrt{2\pi n} (n/e)^n\)を使用すれば
\begin{align*}
  h(\delta) &= \frac{n!}{(np)!(n-np)!} p^{np} (1-p)^{n(1-p)} \exp \left( \frac{n\delta^2}{2p(1-p)} \right) \\
  &\approx \sqrt{\frac{n}{2\pi p(1-p)}} \exp \left( \frac{n\delta^2}{2p(1-p)} \right) .
\end{align*}
したがって
\[ \delta \sim \mathcal{N} \left( 0, \frac{p(1-p)}{n} \right) . \]

\setcounter{chapter}{4}
\chapter{誤差測定のある線形モデル}
\section*{練習問題}
\addcontentsline{toc}{section}{練習問題}
\paragraph{問題5.7}
\(\{\theta_i\}\), \(\{\delta_i\}\), \(\{\epsilon_i\}\)の確率密度函数は
\[
f(\{\theta_i\}, \{\delta_i\}, \{\epsilon_i\})
= \frac{1}{(2\pi)^{3n/2} (\sigma \sigma_x \sigma_y)^n}
\prod_{i=1}^n \exp \left[ - \frac{(\theta_i-\mu)^2}{2\sigma^2} - \frac{\delta_i{}^2}{2\sigma_x{}^2} - \frac{\epsilon_i{}^2}{2\sigma_y{}^2} \right] .
\]
\(\delta_i = x_i - \theta_i\)および\(\epsilon_i = y_i - \alpha - \beta\theta_i\)により変数変換をする．変換のJacobianは
\[
\left\lvert \frac{\partial (\theta_i, \delta_i, \epsilon_i)}{\partial (\theta_i, x_i, y_i)} \right\rvert = 1
\]
なので\(\{\theta_i\}\), \(\{x_i\}\), \(\{y_i\}\)の確率密度函数は
\begin{align*}
  & g(\{\theta_i\}, \{x_i\}, \{y_i\}) \\
  &= \frac{1}{(2\pi)^{3n/2} (\sigma \sigma_x \sigma_y)^n}
  \prod_{i=1}^n \exp \left[ - \frac{(\theta_i-\mu)^2}{2\sigma^2} - \frac{\delta_i{}^2}{2\sigma_x{}^2} - \frac{\epsilon_i{}^2}{2\sigma_y{}^2} \right] \\
  %
  &= \frac{1}{(2\pi)^{3n/2} (\sigma \sigma_x \sigma_y)^n}
  \prod_{i=1}^n \exp \left[
    - \frac{(\theta_i-\mu)^2}{2\sigma^2}
    - \frac{(x_i - \theta_i)^2}{2\sigma_x{}^2}
    - \frac{(y_i - \alpha - \beta\theta_i)^2}{2\sigma_y{}^2} \right] \\
  %
  &= \frac{1}{(2\pi)^{3n/2} (\sigma \sigma_x \sigma_y)^n}
  \prod_{i=1}^n \\
  &\quad\times \exp \left[
    - \left( \frac{1}{\sigma^2} + \frac{1}{\sigma_x{}^2} + \frac{\beta^2}{\sigma_y{}^2} \right) \frac{\theta_i{}^2}{2}
    + \left( \frac{\mu}{\sigma^2} + \frac{x_i}{\sigma_x{}^2} + \frac{\beta(y_i-\alpha)}{\sigma_y{}^2} \right) \theta_i
    - \left( \frac{\mu^2}{2\sigma^2} + \frac{x_i{}^2}{2\sigma_x{}^2} + \frac{(y_i-\alpha)^2}{2\sigma_y{}^2} \right) \right] \\
  %
  &= \frac{1}{(2\pi)^{3n/2} (\sigma \sigma_x \sigma_y)^n}
  \prod_{i=1}^n \exp \left[ - \frac{A}{2} \theta_i{}^2 + B_i \theta_i - C_i \right] .
\end{align*}
これを\(\theta_i\)で積分すれば
\begin{align*}
  g(\{\theta_i\}, \{x_i\}, \{y_i\})
  &= \int d\theta_1 \cdots d\theta_n \,  g(\{\theta_i\}, \{x_i\}, \{y_i\}) \\
  &= \frac{1}{(2\pi)^n (\sigma \sigma_x \sigma_y)^n} \frac{1}{A^{n/2}}
  \prod_{i=1}^n \exp \left( \frac{B_i{}^2}{2A} - C_i \right) .
\end{align*}
よって対数尤度は
\[
l = - n \log(2\pi\sigma_x \sigma_y) - n \log\sigma - \frac{n}{2} \log A + \frac{1}{2A} \sum_i B_i{}^2 - \sum_i C_i .
\]

\(\alpha\), \(\beta\), \(\mu\), \(\sigma\)を最尤推定すればよい．

\chapter{一般化線形モデル}
\setcounter{section}{5}
\section{over-dispersion}
個人間差がある場合は固定効果モデルより混合効果モデルを使用すべき，という話．

\paragraph{(6.42)}
ここでは
\[ E(\mu_i) = \mu , \quad \Var(\mu_i) = \tau^2 , \quad y_i\mid\mu_i \sim \Bin(\mu_i, m_i) \]
を考えている．

\section*{練習問題}
\addcontentsline{toc}{section}{練習問題}
\paragraph{問題6.3}
ここで考えているのはCox-Snell型の疑似決定係数\(R^2_\text{CS}\)．
McFaddenの擬似決定係数
\[ R^2_\text{Mc} = 1- \frac{\ln L(\hat{\boldsymbol\beta})}{\ln L(\boldsymbol{0})} \]
は線形正規モデルにおいては，線形回帰の決定係数\(R^2\)に一致する．

\chapter{ノンパラメトリック回帰モデル}
\setcounter{section}{4}
\section{smootherのバラツキとsmoothingパラメータ}
\paragraph{(7.27)}
\cite{hall1990variance}で導入されたようである．

\setcounter{chapter}{11}
\chapter{トピックスIV}
\setcounter{section}{1}
\section{balanced dataでの推測}
\paragraph{(12.49)(12.50)(12.51)}
まず施設\(i\)のみを考える．
\[
\boldsymbol{1}_n =
\begin{pmatrix}
  1 \\ \vdots \\ 1
\end{pmatrix}
, \quad
I_n =
\begin{pmatrix}
  1 & & \\
  & \ddots & \\
  & & 1
\end{pmatrix}
, \quad
J_n =
\begin{pmatrix}
  1 & \cdots & 1 \\
  \vdots & \ddots & \vdots \\
  1 & \cdots & 1
\end{pmatrix}
\]
とすれば
\[
  \boldsymbol{y}_i
  =
  \begin{pmatrix}
    y_{i11} \\ \vdots \\ y_{iJn}
  \end{pmatrix}
  = (\boldsymbol{1}_J \otimes I_n)
  \begin{pmatrix}
    \beta_1 \\ \vdots \\ \beta_J
  \end{pmatrix}
  + \boldsymbol{1}_{Jn} \alpha_i
  + (\boldsymbol{1}_n \otimes I_J)
  \begin{pmatrix}
    \gamma_{i1} \\ \vdots \\ \gamma_{iJ}
  \end{pmatrix}
  +
  \begin{pmatrix}
    \epsilon_{i11} \\ \vdots \\ \epsilon_{iJn}
  \end{pmatrix}
  \sim \mathcal{N}(X \boldsymbol\beta, V)
\]
である．ただし，
\[
X = I_J \otimes \boldsymbol{1}_n , \quad
\boldsymbol\beta =
\begin{pmatrix}
  \beta_1 \\ \vdots \\ \beta_J
\end{pmatrix}
\]
および
\begin{align*}
  V &= \sigma_\alpha{}^2 \boldsymbol{1}_{Jn} \boldsymbol{1}_{Jn}^\top
  + \sigma_\gamma{}^2 (\boldsymbol{1}_n \otimes I_J) (\boldsymbol{1}_n \otimes I_J)^\top
  + \sigma_E{}^2 I_{Jn} \\
  &= \sigma_\alpha{}^2 J_{Jn} + \sigma_\gamma{}^2 (I_J \otimes J_n) + \sigma_E{}^2 I_{Jn}
\end{align*}
である．

\begin{screen}
  \(V\)の固有値は
  \[
  \lambda_1 = \sigma_E{}^2 , \quad
  \lambda_2 = \sigma_E{}^2 + n\sigma_\gamma{}^2 , \quad
  \lambda_3 = \sigma_E{}^2 + n\sigma_\gamma{}^2 + Jn \sigma_\alpha{}^2 .
  \]
  で与えられ，それぞれに対応する射影行列・重複度は
  \[
    P_1 = I_J \otimes \left( I_n - \frac{1}{n} J_n \right) , \quad
    P_2 = \frac{1}{n} \left( I_J - \frac{1}{J} J_J \right) \otimes J_n , \quad
    P_3 = \frac{1}{Jn} J_{Jn}
  \]
  および
  \[ n_1 = J(n-1) , \quad n_2 = J-1, \quad n_3 = 1 . \]
\end{screen}
\begin{proof}
  固有値を\(\lambda\)，固有ベクトルを
  \[
  \begin{pmatrix}
    x_{11} \\ \vdots \\ x_{Jn}
  \end{pmatrix}
  \]
  とする．固有値方程式は
  \begin{align*}
    \lambda
    \begin{pmatrix}
      x_{11} \\ \vdots \\ x_{Jn}
    \end{pmatrix}
    &= V
    \begin{pmatrix}
      x_{11} \\ \vdots \\ x_{Jn}
    \end{pmatrix}
    \\
    &= \sigma_\alpha{}^2
    J_{Jn}
    \begin{pmatrix}
      x_{11} \\ \vdots \\ x_{Jn}
    \end{pmatrix}
    +
    \sigma_\gamma{}^2
    \begin{pmatrix}
      J_n & & \\
      & \ddots & \\
      & & J_n
    \end{pmatrix}
    \begin{pmatrix}
      x_{11} \\ \vdots \\ x_{Jn}
    \end{pmatrix}
    +
    \sigma_E{}^2
    \begin{pmatrix}
      x_{11} \\ \vdots \\ x_{Jn}
    \end{pmatrix}
    \\
    %
    &= \sigma_\alpha{}^2
    \begin{pmatrix}
      \sum_{jk} x_{jk} \\ \vdots \\ \sum_{jk} x_{jk}
    \end{pmatrix}
    +
    \sigma_\gamma{}^2
    \begin{pmatrix}
      \sum_k x_{1k} \\ \vdots \\ \sum_k x_{2k} \\ \vdots \\ \sum_k x_{Jk} \\ \vdots
    \end{pmatrix}
    +
    \sigma_E{}^2
    \begin{pmatrix}
      x_{11} \\ \vdots \\ x_{Jn}
    \end{pmatrix}
  \end{align*}
  となる．変形すれば
  \[
  (\lambda - \sigma_E{}^2)
  \begin{pmatrix}
    x_{11} \\ \vdots \\ x_{Jn}
  \end{pmatrix}
  = \sigma_\alpha{}^2
  \begin{pmatrix}
    \sum_{jk} x_{jk} \\ \vdots \\ \sum_{jk} x_{jk}
  \end{pmatrix}
  +
  \sigma_\gamma{}^2
  \begin{pmatrix}
    \sum_k x_{1k} \\ \vdots \\ \sum_k x_{2k} \\ \vdots \\ \sum_k x_{Jk} \\ \vdots
  \end{pmatrix}
  .
  \]
  \(n\)個ずつの成分に注目すれば
  \begin{align*}
    & (\lambda - \sigma_E{}^2) x_{11} = \cdots = (\lambda - \sigma_E{}^2) x_{1n} \\
    & \quad \vdots \\
    & (\lambda - \sigma_E{}^2) x_{J1} = \cdots = (\lambda - \sigma_E{}^2) x_{Jn}
  \end{align*}
  が得られる．
  \begin{enumerate}
    \item \(\lambda = \sigma_E{}^2\)の場合．
    これは明らかに固有値方程式を満たす．
    \(\lambda_2 = \sigma_E{}^2 + n \sigma_\gamma{}^2\)とする．
    固有値方程式に代入すれば
    \[ \sum_k x_{jk} = 0 \quad (j=1, \dots, J) . \]
    これらは独立な\(J\)個の方程式なので部分空間の次元は\(Jn - J = J(n-1)\)である．ここで
    \[
    \boldsymbol{w}_1^{(l)}
    = \frac{1}{\sqrt{l(l+1)}}
    \begin{pmatrix}
      1 \\ \vdots \\ 1 \\ -l \\ 0 \\ \vdots \\ 0
    \end{pmatrix}
    \in \mathbb{R}^n
      \quad (l = 1, \dots, n-1)
    \]
    とする．これらは正規直交系を成す．もとの\(Jn\)次元での正規直交系としては\(\dim = J(n-1)\)の
    \[
      \boldsymbol{v}_1^{(1, l)} =
      \begin{pmatrix}
        \boldsymbol{w}_1^{(l)} \\ 0 \\ \vdots \\ 0
      \end{pmatrix}
      , \quad
      \boldsymbol{v}_1^{(2, l)} =
      \begin{pmatrix}
        0 \\ \boldsymbol{w}_1^{(l)} \\ 0 \\ \vdots \\ 0
      \end{pmatrix}
      , \quad
      \dots
      , \quad
      \boldsymbol{v}_1^{(J, l)} =
      \begin{pmatrix}
        0 \\ \vdots \\ 0 \\ \boldsymbol{w}_1^{(l)}
      \end{pmatrix}
    \]
    が取れる．よって\(\lambda_1\)の対応する射影行列は
    \begin{align*}
      P_1 &= \sum_j \sum_l \boldsymbol{v}_1^{(j, l)} (\boldsymbol{v}_1^{(j, l)} )^\top \\
      &=
      \begin{pmatrix}
        \sum_l \boldsymbol{w}_1^{(l)} (\boldsymbol{w}_1^{(l)})^\top & & \\
        & \ddots & \\
        & & \sum_l \boldsymbol{w}_1^{(l)} (\boldsymbol{w}_1^{(l)})^\top
      \end{pmatrix}
      \\
      &= I_J \otimes \left( I_n - \frac{1}{n} J_n \right) .
    \end{align*}
    最後の計算は帰納法により容易に示せる．

    \item \(\lambda \neq \sigma_E{}^2\)の場合．
    \begin{align}
      x_{11} = \cdots = x_{1n} , \quad \ldots , \quad x_{J1} = \cdots = x_{Jn}
      \label{12_49_p2p3}
    \end{align}
    となる．これらを再び固有値方程式に代入すれば
    \[
    (\lambda - \sigma_E{}^2)
    \begin{pmatrix}
      x_{11} \\ \vdots \\ x_{Jn}
    \end{pmatrix}
    = n \sigma_\alpha{}^2
    \begin{pmatrix}
      \sum_{j} x_{j1} \\ \vdots \\ \sum_{j} x_{j1}
    \end{pmatrix}
    +
    n \sigma_\gamma{}^2
    \begin{pmatrix}
      x_{11} \\ \vdots \\ x_{21} \\ \vdots \\ x_{J1} \\ \vdots
    \end{pmatrix}
    .
    \]
    整理すれば
    \[
    (\lambda - \sigma_E{}^2 - n \sigma_\gamma{}^2)
    \begin{pmatrix}
      x_{11} \\ x_{21} \\ \vdots \\ x_{J1}
    \end{pmatrix}
    =  n \sigma_\alpha{}^2
    \begin{pmatrix}
      \sum_{j} x_{j1} \\ \vdots \\ \sum_{j} x_{j1}
    \end{pmatrix}
    .
    \]
    右辺の成分はすべて同じなので
    \[
    (\lambda - \sigma_E{}^2 - n \sigma_\gamma{}^2) x_{11} = \cdots = (\lambda - \sigma_E{}^2 - n \sigma_\gamma{}^2) x_{J1} .
    \]
    \begin{enumerate}
      \item \(\lambda = \sigma_E{}^2 + n \sigma_\gamma{}^2\)の場合．
      これは明らかに固有値方程式を満たす．
      \(\lambda_2 = \sigma_E{}^2 + n \sigma_\gamma{}^2\)とする．
      固有値方程式に代入すれば
      \[ \sum_j x_{j1} = 0 . \]
      これと\eqref{12_49_p2p3}を併せれば，固有値ベクトルは\(k=1\)の成分のみで完全に決定する．
      直交系としては
      \[
      \begin{pmatrix}
        x_{11} \\ \vdots \\ x_{J1}
      \end{pmatrix}
      = \boldsymbol{w}_2^{(l)}
      = \begin{pmatrix}
        1 \\ \vdots \\ 1 \\ -l \\ 0 \\ \vdots \\ 0
      \end{pmatrix}
       \quad (l = 1, \dots, J-1)
      \]
      が取れる．よって，もとの\(Jn\)次元での正規直交系としては\(\dim = J-1\)の
      \[
      \begin{pmatrix}
        x_{11} \\ \vdots \\ x_{Jn}
      \end{pmatrix}
      = \boldsymbol{v}_2^{(l)}
      = \frac{1}{\sqrt{l(l+1)n}} \boldsymbol{w}_2^{(l)} \otimes \boldsymbol{1}_n
      \quad (k = 1, \dots, J-1)
      \]
      が取れる．よって\(\lambda_2\)に対応する射影行列は
      \begin{align*}
        P_2
        &= \sum_{l=1}^{J-1} \boldsymbol{v}_2^{(l)} (\boldsymbol{v}_2^{(l)})^\top
        = \sum_{l=1}^{J-1} \frac{1}{l(l+1)n} (\boldsymbol{w}_2^{(l)} \otimes \boldsymbol{1}_n) (\boldsymbol{w}_2^{(l)} \otimes \boldsymbol{1}_n)^\top \\
        %
        &= \sum_{l=1}^{J-1} \frac{1}{l(l+1)n} (\boldsymbol{w}_2^{(l)} \otimes \boldsymbol{1}_n) ((\boldsymbol{w}_2^{(l)})^\top \otimes \boldsymbol{1}_n^\top) \\
        %
        &= \frac{1}{n} \sum_{l=1}^{J-1} \frac{1}{l(l+1)} (\boldsymbol{w}_2^{(l)} (\boldsymbol{w}_2^{(l)})^\top) \otimes (\boldsymbol{1}_n \boldsymbol{1}_n^\top) \\
        %
        &= \frac{1}{n} \left( I_J - \frac{1}{J} J_J \right) \otimes J_n .
      \end{align*}
      最後の計算は帰納法により容易に示せる．

      \item \(\lambda \neq \sigma_E{}^2 + n \sigma_\gamma{}^2\)の場合．
      \(x_{11} = \cdots = x_{J1} =\)である．\eqref{12_49_p2p3}と併せれば
      \[ x_{11} = \cdots x_{1n} = x_{21} = \cdots = x_{Jn} . \]
      固有ベクトルは
      \[ \boldsymbol{v}_3 = \frac{1}{\sqrt{Jn}} \boldsymbol{1}_{Jn} \]
      の1本のみ．固有値方程式に代入すれば固有値は
      \[ \lambda_3 =  \sigma_E{}^2 + n \sigma_\gamma{}^2 + Jn \sigma_\alpha{}^2 . \]
      対応する射影行列は
      \[ P_3 = \boldsymbol{v}_3 \boldsymbol{v}_3^\top = \frac{1}{Jn} J_{Jn} . \]
    \end{enumerate}
  \end{enumerate}
\end{proof}

射影行列の性質から
\[ P_1 + P_2 + P_3 = I_{Jn} , \quad P_aP_a = P_a , \quad P_aP_b = 0 \quad (a \neq b) \]
を満たす．よって
\begin{align}
  \det V = \lambda_1{}^{n_1} \lambda_2{}^{n_2} \lambda_3{}^{n_3} , \quad
  V^{-1} = \frac{1}{\lambda_1} P_1 + \frac{1}{\lambda_2} P_2 + \frac{1}{\lambda_3} P_3
  \label{12_49_V_decomp_corollary}
\end{align}
が得られる．

施設\(i\)の対数尤度は
\[
l_i = -\frac{Jn}{2} \log 2\pi - \frac{1}{2} \log \lvert V \rvert
 - \frac{1}{2} (\boldsymbol{y}_i - X \boldsymbol\beta)^\top V^{-1} (\boldsymbol{y}_i - X \boldsymbol\beta) .
\]
\(\boldsymbol{y}_i\)は独立なので，全施設の対数尤度は
\begin{align}
  l = -\frac{IJn}{2} \log 2\pi - \frac{I}{2} \log \lvert V \rvert
    - \frac{1}{2} \sum_{i=1}^I (\boldsymbol{y}_i - X \boldsymbol\beta)^\top V^{-1} (\boldsymbol{y}_i - X \boldsymbol\beta) .
  \label{12_49_ll_all_inst}
\end{align}
まず\(\boldsymbol\beta\)の最尤推定量を求める．
\[
\frac{\partial l}{\partial \beta} = \sum_i X^\top V^{-1} (\boldsymbol{y}_i - X \boldsymbol\beta) = \boldsymbol{0}
\]
なので，
\[
\hat{\boldsymbol\beta} = (X^\top V^{-1}X)^{-1} X^\top V^{-1} \bar{\boldsymbol{y}}_\bullet
, \quad
\bar{\boldsymbol{y}}_\bullet := \frac{1}{I} \sum_i \boldsymbol{y}_i .
\]
ここで
\begin{align*}
  X^\top V^{-1}
  &= (I_J \otimes \boldsymbol{1}_n^\top)
  \left( \frac{1}{\lambda_1} P_1 + \frac{1}{\lambda_2} P_2 + \frac{1}{\lambda_3} P_3 \right) \\
  &= \frac{1}{\lambda_1} (I_J \otimes \boldsymbol{1}_n^\top) \left( I_J \otimes \left( I_n - \frac{1}{n} J_n \right) \right)
  + \frac{1}{n\lambda_2} (I_J \otimes \boldsymbol{1}_n^\top) \left( \left( I_J - \frac{1}{J} J_J \right) \otimes J_n \right) \\
  &\quad + \frac{1}{Jn\lambda_3} (I_J \otimes \boldsymbol{1}_n^\top) \left( J_J \otimes J_n \right) \\
  %
  &= \frac{1}{\lambda_1} I_J \otimes \boldsymbol{1}_n^\top \left( I_n - \frac{1}{n} J_n \right)
  + \frac{1}{n\lambda_2} \left( I_J - \frac{1}{J} J_J \right) \otimes \boldsymbol{1}_n^\top J_n
  + \frac{1}{Jn\lambda_3} J_J \otimes \boldsymbol{1}_n^\top J_n \\
  %
  &= \frac{1}{\lambda_2} \left( I_J - \frac{1}{J} J_J \right) \otimes \boldsymbol{1}_n^\top
  + \frac{1}{J\lambda_3} J_J \otimes \boldsymbol{1}_n^\top .
\end{align*}
さらに
\begin{align*}
  X^\top V^{-1} X &=
  \left[ \frac{1}{\lambda_2} \left( I_J - \frac{1}{J} J_J \right) \otimes \boldsymbol{1}_n^\top
  + \frac{1}{J\lambda_3} J_J \otimes \boldsymbol{1}_n^\top \right] \otimes (I_J \otimes \boldsymbol{1}_n) \\
  %
  &= \frac{1}{\lambda_2} \left( I_J - \frac{1}{J} J_J \right) \otimes n
  + \frac{1}{J\lambda_3} J_J \otimes n \\
  %
  &= \frac{n}{\lambda_2} \left( I_J - \frac{1}{J} J_J \right) + \frac{n}{J\lambda_3} J_J .
\end{align*}
以上から
\[ X^\top V^{-1} = \frac{1}{n} (X^\top V^{-1} X ) \otimes \boldsymbol{1}_n^\top . \]
よって
\begin{align*}
  \hat{\boldsymbol\beta} &= (X^\top V^{-1}X)^{-1} X^\top V^{-1} \bar{\boldsymbol{y}}_\bullet \\
  &= \frac{1}{n}(X^\top V^{-1}X)^{-1} \left[ (X^\top V^{-1} X ) \otimes \boldsymbol{1}_n^\top \right] \bar{\boldsymbol{y}}_\bullet \\
  %
  &= \frac{1}{n} [(X^\top V^{-1}X)^{-1} \otimes 1] \left[ (X^\top V^{-1} X ) \otimes \boldsymbol{1}_n^\top \right] \bar{\boldsymbol{y}}_\bullet \\
  &= \frac{1}{n} I_J \otimes \boldsymbol{1}_n^\top \bar{\boldsymbol{y}}_\bullet \\
  %
  &= \frac{1}{n}
  \begin{pmatrix}
    1 & \cdots & 1 &&& & \\
    &&& \ddots &&& \\
    &&&& 1 & \cdots & 1
  \end{pmatrix}
  \begin{pmatrix}
    \bar{y}_{\bullet11} \\ \vdots \\ \bar{y}_{\bullet Jn}
  \end{pmatrix}
  \\
  &=
  \begin{pmatrix}
    \bar{y}_{\bullet 1 \bullet} \\ \vdots \\ \bar{y}_{\bullet J \bullet}
  \end{pmatrix} .
\end{align*}
従って
\[
\boldsymbol{y}_i - X \hat{\boldsymbol\beta}
=
\begin{pmatrix}
  y_{i11} - \bar{y}_{\bullet 1 \bullet} \\
  \vdots \\
  y_{i1n} - \bar{y}_{\bullet 1 \bullet} \\ \hline
  y_{i21} - \bar{y}_{\bullet 2 \bullet} \\
  \vdots \\
  y_{i2n} - \bar{y}_{\bullet 2 \bullet} \\ \hline
  \vdots \\
  y_{iJn} - \bar{y}_{\bullet J \bullet}
\end{pmatrix}
=:
\begin{pmatrix}
  \boldsymbol{z}_{i1} \\ \vdots \\  \boldsymbol{z}_{iJ}
\end{pmatrix}
.
\]
\(\boldsymbol{z}_{ij}\)の第\(k\)成分は
\[ (\boldsymbol{z}_{ij})_k = y_{ijk} - \bar{y}_{\bullet j \bullet} \]
である．
\(P_1\)との積は
\[
  P_1 (\boldsymbol{y}_i - X \hat{\boldsymbol\beta})
  =
  \begin{pmatrix}
    I_n - J_n/n & & \\
    & \ddots & \\
    & & I_n - J_n/n
  \end{pmatrix}
  \begin{pmatrix}
    \boldsymbol{z}_{i1} \\ \vdots \\  \boldsymbol{z}_{iJ}
  \end{pmatrix}
  =
  \begin{pmatrix}
    (I_n - J_n/n) \boldsymbol{z}_{i1} \\
    \vdots \\
    (I_n - J_n/n) \boldsymbol{z}_{iJ}
  \end{pmatrix}
  .
\]
\(j\)番目のblockは
\[ \left( I_n - \frac{J_n}{n} \right) \boldsymbol{z}_{ij} \]
であり，これの第\(k\)成分は
\[
(\boldsymbol{z}_{ij})_k - \frac{1}{n} \sum_k (\boldsymbol{z}_{ij})_k
= (y_{ijk} - \bar{y}_{\bullet j \bullet}) - \frac{1}{n} \sum_k (y_{ijk} - \bar{y}_{\bullet j \bullet})
= y_{ijk} - \bar{y}_{ij\bullet} .
\]
よって，
\begin{align}
  \begin{split}
    \sum_i (\boldsymbol{y}_i - X \hat{\boldsymbol\beta})^\top P_1 (\boldsymbol{y}_i - X \hat{\boldsymbol\beta})
    &= \sum_i(\boldsymbol{y}_i - X \hat{\boldsymbol\beta})^\top P_1 P_1 (\boldsymbol{y}_i - X \hat{\boldsymbol\beta}) \\
    &= \sum_i \sum_{jk} (y_{ijk} - \bar{y}_{ij\bullet} )^2 \\
    &= \text{SSE} .
  \end{split}
  \label{12_49_p1_sse}
\end{align}
次に\(P_2\)との積は
\[
  P_2 (\boldsymbol{y}_i - X \hat{\boldsymbol\beta})
  = \frac{1}{n}
  \begin{pmatrix}
    (1-1/J) J_n & - J_n/J & \cdots \\
    - J_n/J & (1-1/J) J_n & \cdots \\
    \vdots & \vdots & \ddots
  \end{pmatrix}
  \begin{pmatrix}
    \boldsymbol{z}_{i1} \\ \vdots \\  \boldsymbol{z}_{iJ}
  \end{pmatrix}
  .
\]
これの\(j\)番目のblockは
\[ \frac{1}{n} J_n \boldsymbol{z}_{ij} - \frac{1}{Jn} \sum_j J_n \boldsymbol{z}_{ij} \]
であり，これの第\(k\)成分は
\[
\frac{1}{n} \sum_k (\boldsymbol{z}_{ij})_k - \frac{1}{Jn} \sum_{jk} (\boldsymbol{z}_{ij})_k
= \frac{1}{n} \sum_k (y_{ijk} - \bar{y}_{\bullet j\bullet}) - \frac{1}{Jn} \sum_{jk} (y_{ijk} - \bar{y}_{\bullet j\bullet})
= \bar{y}_{ij\bullet} - \bar{y}_{\bullet j\bullet} - \bar{y}_{i\bullet\bullet} + \bar{y}_{\bullet\bullet\bullet} .
\]
よって
\begin{align}
  \begin{split}
    \sum_i (\boldsymbol{y}_i - X \hat{\boldsymbol\beta})^\top P_2 (\boldsymbol{y}_i - X \hat{\boldsymbol\beta})
    &= \sum_i(\boldsymbol{y}_i - X \hat{\boldsymbol\beta})^\top P_2 P_2 (\boldsymbol{y}_i - X \hat{\boldsymbol\beta}) \\
    &= \sum_i \sum_{jk} (\bar{y}_{ij\bullet} - \bar{y}_{\bullet j\bullet} - \bar{y}_{i\bullet\bullet} + \bar{y}_{\bullet\bullet\bullet})^2 \\
    &= \text{SSAB} .
  \end{split}
  \label{12_49_p2_ssab}
\end{align}

最後に\(P_3\)との積は
\[
  P_3 (\boldsymbol{y}_i - X \hat{\boldsymbol\beta})
  = \frac{1}{Jn}
  \begin{pmatrix}
    1 & \cdots & 1 \\
    \vdots & \ddots & \vdots \\
    1 & \cdots & 1
  \end{pmatrix}
  \begin{pmatrix}
    \boldsymbol{z}_{i1} \\ \vdots \\  \boldsymbol{z}_{iJ}
  \end{pmatrix}
  .
\]
これの\((j-1)n+k\)成分は
\[
\frac{1}{Jn} \sum_{jk} (\boldsymbol{z}_{ij})_k
= \frac{1}{Jn} \sum_{jk} (y_{ijk} - \bar{y}_{\bullet j \bullet})
= \bar{y}_{i\bullet\bullet} - \bar{y}_{\bullet\bullet\bullet} .
\]
よって
\begin{align}
  \begin{split}
    \sum_i (\boldsymbol{y}_i - X \hat{\boldsymbol\beta})^\top P_3 (\boldsymbol{y}_i - X \hat{\boldsymbol\beta})
    &= \sum_i(\boldsymbol{y}_i - X \hat{\boldsymbol\beta})^\top P_3 P_3 (\boldsymbol{y}_i - X \hat{\boldsymbol\beta}) \\
    &= \sum_i \sum_{jk} (\bar{y}_{i\bullet\bullet} - \bar{y}_{\bullet\bullet\bullet})^2 \\
    &= \text{SSA} .
  \end{split}
  \label{12_49_p3_ssa}
\end{align}

\eqref{12_49_V_decomp_corollary}\eqref{12_49_ll_all_inst}\eqref{12_49_p1_sse}\eqref{12_49_p2_ssab}\eqref{12_49_p3_ssa}から
\begin{align*}
  l &= - \frac{IJn}{2} \log (2\pi)
  - \frac{I}{2} \log (\lambda_1{}^{n_1} \lambda_2{}^{n_2} \lambda_3{}^{n_3})
  - \frac{1}{2} \sum_{\nu=1}^3 \lambda_\nu \sum_i (\boldsymbol{y}_i - X \hat{\boldsymbol\beta})^\top P_\nu (\boldsymbol{y}_i - X \hat{\boldsymbol\beta}) \\
  &=  - \frac{IJn}{2} \log (2\pi)
  - \frac{IJ(n-1)}{2} \log\lambda_1 - \frac{I(J-1)}{2} \log\lambda_2 - \frac{I}{2} \log\lambda_3
  - \frac{1}{2} \left( \frac{\text{SSE}}{\lambda_1} + \frac{\text{SSAB}}{\lambda_2} + \frac{\text{SSA}}{\lambda_3} \right) .
\end{align*}
\(\lambda_1\), \(\lambda_2\), \(\lambda_3\)の最尤推定量は
\[
{\hat\lambda}_1 = \text{MSE} , \quad
{\hat\lambda}_2 = \frac{I-1}{I} \text{MSAB} , \quad
{\hat\lambda}_3 = \frac{I-1}{I} \text{MSA} .
\]
よって
\begin{align*}
  {\hat\sigma}_E{}^2 &= {\hat\lambda}_1 = \text{MSE} , \\
  {\hat\sigma}_\gamma{}^2 &= \frac{{\hat\lambda}_2 - {\hat\lambda}_1}{n} = \frac{1}{n} \left( \frac{I-1}{I} \text{MSAB} - \text{MSE} \right) , \\
  {\hat\sigma}_\alpha{}^2 &= \frac{{\hat\lambda}_3 - {\hat\lambda}_2}{n} = \frac{I-1}{IJn} (\text{MSA} - \text{MSAB}) .
\end{align*}
